\documentclass[smallextended]{svjour3}       % onecolumn (second format)

\usepackage[a-1b]{pdfx}

\smartqed  % flush right qed marks, e.g. at end of proof

\usepackage{graphicx}
%\graphicspath{{.},{../}}

\usepackage{mathptmx}      % use Times fonts if available on your TeX system

% insert here the call for the packages your document requires

\usepackage{amssymb}
\usepackage{mathtools}
\usepackage{mathrsfs}
\usepackage{rsfso}

%\usepackage{amsthm}

\usepackage{xcolor}
\usepackage{tikz}

\renewcommand{\div}{\nabla\cdot}
\renewcommand{\Re}{{\mathbb R}}
\newcommand{\bfx}{{\mathbf x}}

\newcommand*\Laplace{\mathop{}\!\mathbin\bigtriangleup} 

\definecolor{myBlue}{RGB}{190,203,211}
\definecolor{myRed}{RGB}{244,186,175}
\definecolor{myGreen}{RGB}{114,138,122}
\definecolor{myYellow}{RGB}{241,208,157}

\begin{document}

\section{Problem setting}
A scalar advection-diffusion problem without reaction.
\begin{equation} \label{eq:ad}
		  u_t + \div (F(u) - D(u) \nabla u) = 0, \ \bfx\in\Re^d
\end{equation}
The variable $D(u)$ will be treated with  Kirchhoff transformation.
The diffusion coefficient $D$ is treated as a constant here.
This will not affect the following discussion on numerical scheme.
\begin{equation} \label{eq:sad}
		  u_t + \div F(u) - D\Laplace u = 0, \ \bfx\in\Re^d
\end{equation}
\subsection{Finite volume scheme}
Replace scalar $u$ with cell averaged value 
\begin{equation}\label{eq:cellaveu}
		  \bar u = \frac{1}{|E|} \int_E u(\bfx, t) d\bfx
\end{equation}
and Lax-Friedrichs numerical flux
\begin{equation}\label{eq:lax}
		  f(u) = \frac{1}{2} [F(u^+)  + F(u^-) \cdot \nu_E - \alpha (u^+ - u^-)]
\end{equation}
The finite volume scheme of advection-diffusion problem is:
\begin{equation}\label{eq:fv}
		  \bar u_t + \frac{1}{|E|} \int_{\partial E} (\hat f(\bar u) - \nabla \bar u \cdot \nu_E) d\sigma(\bfx) = 0
\end{equation}
In 2D quadrilateral logically rectangular mesh, 
we apply gaussian quadrature every edge as the approximation of integral of flux. 
I use Multi-level WENO reconstruction scheme to reconstruct the point value in 
advective numerical flux.
\begin{equation}\label{eq:laxWENO}
		  \hat f(u) = \frac{1}{2} [F(\mathscr{R}^+(\{\bar u\}_+)) + F(\mathscr{R}^-(\{\bar u\}_-))\cdot \nu_E - 
		  \alpha(\mathscr{R}^+(\{\bar u\}_+) - \mathscr{R}^-(\{\bar u\}_-))] 
\end{equation}
Let $\{\bar u\}_+ = \{\bar u_i \mid i \in \text{WENO stencil} \quad I_+ \}$. 
Diffusive flux is different from its advective counterpart.
I reconstruct derivative across the edge by sampling along the normal direction.

\section{Computation of Jacobian}
Implicit method contains solving a non linear system and a linear subproblem.
I write the general non linear system as:
\begin{equation} \label{eq:nonlinear}
		  \bar u_t + \mathscr{F}(\bar u) - \mathscr{G}(\bar u) = 0
\end{equation}
where $\mathscr{F}(\bar u)$ represents advection part, and $\mathscr{G} (\bar u)$ represents diffusion part.
The semi discretized form is 
\begin{equation}
		  \bar u^{n+1} - \bar u^{n} + \Delta t [\mathscr{F} (\bar u^{n+1}) - \mathscr{G}(\bar u^{n+1})] = 0
\end{equation}
There are many well-researched numerical methods that can be used to solve such non linear system.
%I use non linear solver provided by PETSc.
%I provide routine computing right hand side and jacobian matrix.
%PETSc provides line search and preconditioner algorithm.
\subsection{ML-WENO reconstruction}
I use ML-WENO reconstruction scheme to reconstruct point values
from cell-averaged value...
I use two stage non linear weighting. The first stage is
calculated as
\begin{equation}\label{eq:fstage}
		  \check \omega_l = \frac{\omega_l}{(\sigma_l + \epsilon h^2_0)^{\eta_l}} 
		  \quad \eta_l = \left \{\begin{aligned}
					 &1, \quad \text{if } r=1\\
					 &3, \quad \text{if } r=2\\
					 &4, \quad \text{if } r\geq 3
		  \end{aligned}
		  \right .
\end{equation}
\begin{equation}\label{eq:fstagenorm}
		  \bar \omega_l = \frac{\check \omega_l}{\sum_k \check \omega_k}
\end{equation}
The the second stage is calculated based on first stage:
\begin{equation}\label{eq:sstage}
		  \hat \omega_l = \frac{\bar \omega_l}{(\sigma_l + \epsilon h^2_0)^{sr_l}}
\end{equation}
\begin{equation}\label{eq:sstagenorm}
		  \tilde \omega_l = \frac{\hat \omega_l}{\sum_k \hat \omega_k}
\end{equation}

I write the reconstruction as combination of polynomials:
\begin{equation}\label{eq:weno}
		  \mathscr{R}(\bfx) = \sum_{l\in L} P_l(\bfx) \tilde\omega_l(\bar u)
\end{equation}
Suppose the reconstruction happens in a given stencil $I$. 
The cell averaged values $\bar u$ in this stencil is represented as $\{\bar u_s\}_{s \in S}$.
The derivative of this reconstruction with respect to $\bar u_s$ is:
\begin{equation}
		  \frac{\partial \mathscr{R}}{\partial \bar u_s} = 
		  \sum_{l \in L} \Big(\frac{\partial P_l}{\partial \bar u_s}\tilde \omega_l + P_l\frac{\partial \tilde \omega_l}{\partial \bar u_s}\Big) 
\end{equation}
It consists of two parts, derivative of polynomials $\frac{\partial P_l}{\partial \bar u_s}$ and 
derivative of non linear weights $\frac{\partial \tilde \omega_l}{\partial \bar u_s}$.
The stencil polynomial is a combination of single polynomials satisfying 
\begin{equation}\label{eq:spolyn}
		  \frac{1}{|E_k|} \int_{E_k} p_s(\bfx) d\bfx = \left \{ 
		  \begin{aligned}
					 &1 \quad s=k\\
					 &0 \quad s\neq k
		  \end{aligned}\right .
\end{equation}
The stencil polynomial is expressed in terms of combination of all these polynomials
\begin{equation}\label{eq:Spolyn}
		  P(\bfx) = \sum_{s\in S} p_s(\bfx) \bar u_s
\end{equation}
The derivative of polynomial is trivial. The derivative of non linear weight is 
a little more complicated.
\begin{equation}\label{eq:derivative}
		  \frac{\partial \mathscr{R}}{\partial \bar u_s} = \sum_{l \in L} p_s(\bfx) \tilde \omega_l + 
		  \sum_{l\in L} P_l\frac{\partial \tilde \omega_l}{\partial \bar u_s}
\end{equation}


\section{Implicit time stepping}

\subsection{Backward Euler}

\subsection{SSP time stepping}

\section{Parallelism}
Parallel solver in inevitable in large scale computation.
Figure~\ref{fig:fullmesh} 
shows a full view of quadrilateral (logically rectangular) mesh.
\begin{figure}[htb]
		  \parbox{\textwidth}{\begin{minipage}[t]{0.45\textwidth}
		  \input mesh.tex
		  \caption{Logically rectangular 
		  mesh displayed in full. The index of 
		  a stencil (colored cells) is the 
		  same as the index of the red cell.}
\label{fig:fullmesh}
		  \end{minipage}
		  \hfill
		  \begin{minipage}[t]{.45\textwidth}
\input stencil.tex
					 \caption{Two sample stencils used by cell $E_{i,j}$.
					 The index of these two stencils are $S_{i-2,j}$ and $S_{i,j-2}$.
					 Consistent with the index of bottom left cell.}
		  \label{fig:stencil}
		  \end{minipage}}
\end{figure}
Counting all stencils without repeating is convenient 
in the logically rectangular mesh setting. 
We use the index of bottom left cell as the index of 
the corresponding stencil. A $\text{M}\times \text{N}$ mesh will then contain
$(\text{M}-\text{I}+1)\times(\text{N}-\text{J}+1)$ stencils 
of size $\text{I}\times \text{J}$.

\subsection{Mesh partition}
I examine the following partition of mesh without lossing generality.
The dashed line indicates ghost cells distributed by
MPI communication. 
The size of ghost layer is determined by the largest stencil
size required for cells in the ghost layer.
Increasing the order of reconstruction (increasing the size
of corresponding stencil) will thus increase the communication 
cost. However, the distribution of mesh only occur once 
in the beginning of the computation. I will not worry about 
it considering we are running on the 2D mesh right now.
\begin{figure}\label{fig:partitionmesh}
\input quadmesh.tex
		  \caption{Illustration of mesh partition.
		  Ghost cells are drawn with dashed lines.
		  The size of ghost layer is determined 
		  by the size of stencils needed by cell in 
		  the ghost layer.}
\end{figure}


\subsection{Sparsity pattern of Jacobian matrix}

\subsection{Optimal preconditioner}

\begin{itemize}
					 \item 1. Incomplete LU decomposition
					 \item 2. Additive schwarz
\end{itemize}

\subsection{Scalability}

\end{document}
